\documentclass[aspectratio=169,10pt]{beamer}

\usetheme{metropolis}
\metroset{sectionpage=none}
\usepackage{amsmath,amssymb,bm,mathtools}
\usepackage{booktabs}
\usepackage{array}
\usepackage{multicol}
\usepackage{hyperref}

\title{Fast KRR via Frequency-Space Preconditioning}
\subtitle{EigenPro-style deflation for Equispaced Fourier Gaussian Processes}
\author{[Your Name]}
\institute{Final Thesis Presentation}
\date{[Date]}

\newcommand{\R}{\mathbb{R}}
\newcommand{\C}{\mathbb{C}}
\newcommand{\E}{\mathbb{E}}
\newcommand{\norm}[1]{\left\lVert #1 \right\rVert}
\newcommand{\kmat}{K}
\newcommand{\Ph}{\Phi}
\newcommand{\sig}{\sigma^2}
\newcommand{\lam}{\lambda}

\begin{document}

% ---------------------------------------------------------
\begin{frame}
  \titlepage
\end{frame}

% ---------------------------------------------------------
\begin{frame}{Thesis message in one slide}
\begin{block}{Problem}
Kernel ridge regression (KRR) is statistically attractive, but standard solvers scale poorly because they require dense kernel linear algebra.
\end{block}

\begin{block}{Existing progress}
\begin{itemize}
  \item \textbf{EigenPro}: accelerates iterative kernel solvers by deflating the largest spectral modes.
  \item \textbf{EFGP}: uses Fourier grids and NUFFT/FFT to reduce low-dimensional translation-invariant kernel regression to fast weight-space linear algebra.
\end{itemize}
\end{block}

\begin{block}{My direction}
Put the deflation idea directly in the \textbf{frequency / weight-space system} of EFGP, aiming to keep EFGP's near-linear data pass while reducing CG iteration counts.
\end{block}
\end{frame}

% ---------------------------------------------------------
\begin{frame}{Roadmap for a 30-minute talk}
\begin{enumerate}
  \item KRR fundamentals and why kernel methods are still relevant.
  \item Computational bottleneck: direct solve vs. iterative solve.
  \item EigenPro: spectral deflation as preconditioning.
  \item Translation-invariant kernels and EFGP.
  \item Proposed algorithm: frequency-space deflation preconditioner.
  \item Experiments: speed, conditioning, and the top-$q$ trade-off.
\end{enumerate}
\end{frame}

% ---------------------------------------------------------
\section{Motivation and KRR fundamentals}

\begin{frame}{Why KRR is still useful now}
\begin{itemize}
  \item \textbf{Strong inductive bias}: kernels encode smoothness, locality, periodicity, and invariances explicitly.
  \item \textbf{Theory}: convex objective, unique solution, and a mature generalization theory.
  \item \textbf{Connections to modern ML}: wide neural networks and NTK regimes renew interest in kernel viewpoints.
  \item \textbf{Recent scalability}: Fourier/NUFFT-based methods show that some kernel estimators can reach $O(n\log n)$ time and memory.
\end{itemize}

\vspace{0.5em}
\begin{block}{Key tension}
The statistics are clean, but the computational linear algebra is the bottleneck.
\end{block}
\end{frame}

% ---------------------------------------------------------
\begin{frame}{KRR problem statement}
Given data $(x_i,y_i)_{i=1}^N$, kernel $k$, and RKHS $\mathcal{H}$:
\[
\hat f = \arg\min_{f\in\mathcal{H}}
\frac{1}{N}\sum_{i=1}^{N}(f(x_i)-y_i)^2 + \lambda\norm{f}_{\mathcal{H}}^2.
\]

By the representer theorem,
\[
\hat f(x)=\sum_{i=1}^N \alpha_i k(x,x_i).
\]

The coefficients solve
\[
(K+N\lambda I)\alpha = y,
\quad K_{ij}=k(x_i,x_j).
\]

\begin{block}{Same linear algebra as GP posterior mean}
In GP notation, the posterior mean solve is $(K+\sigma^2 I)\alpha=y$.
\end{block}
\end{frame}

% ---------------------------------------------------------
\begin{frame}{The two computational obstacles}
\begin{columns}[T,totalwidth=\textwidth]
\begin{column}{0.48\textwidth}
\textbf{Direct method}
\[
(K+\sig I)\alpha=y
\]
\begin{itemize}
  \item Store $K$: $O(N^2)$ memory.
  \item Cholesky solve: $O(N^3)$ time.
  \item Practical only for moderate $N$.
\end{itemize}
\end{column}
\begin{column}{0.48\textwidth}
\textbf{Iterative method}
\[
\alpha_{t+1}=\alpha_t+\text{correction}
\]
\begin{itemize}
  \item Avoids explicit inverse.
  \item Needs repeated matrix-vector products.
  \item Dense kernel matvec is $O(N^2)$.
  \item Slow if the system is ill-conditioned.
\end{itemize}
\end{column}
\end{columns}

\vspace{0.6em}
\begin{alertblock}{Target}
Move from a dense $O(N^2)$ bottleneck toward $O(N\log N)$ or $O(N+M\log M)$ workflows.
\end{alertblock}
\end{frame}

% ---------------------------------------------------------
\begin{frame}{Conjugate gradient and condition number}
For a symmetric positive definite system $A z=b$, CG satisfies
\[
\frac{\norm{e_t}_A}{\norm{e_0}_A}
\leq
2\left(\frac{\sqrt{\kappa(A)}-1}{\sqrt{\kappa(A)}+1}\right)^t,
\quad
\kappa(A)=\frac{\lambda_{\max}(A)}{\lambda_{\min}(A)}.
\]

\begin{block}{Interpretation}
The number of iterations scales roughly like
\[
 t = O\!\left(\sqrt{\kappa(A)}\log\frac{1}{\varepsilon}\right).
\]
\end{block}

\begin{alertblock}{For kernel systems}
The largest eigenvalue often grows with data size, while the smallest eigenvalue is controlled by regularization/noise. A typical upper scale is
\[
\kappa(K+\sig I) \lesssim 1+\frac{N}{\sig}.
\]
\end{alertblock}
\end{frame}

% ---------------------------------------------------------
\begin{frame}{Why the bottleneck is not only the matvec}
Even with a fast approximate matvec, the solve can still be slow:
\[
\text{total time} \approx
\underbrace{T_{\rm setup}}_{\text{kernel/feature preprocessing}}+
\underbrace{n_{\rm iter}}_{\text{depends on }\kappa}
\times
\underbrace{T_{\rm matvec}}_{\text{dense, FFT, or NUFFT}}.
\]

\begin{itemize}
  \item Classical CG with dense kernel: $T_{\rm matvec}=O(N^2)$.
  \item EFGP: $T_{\rm matvec}=O(M\log M)$, independent of $N$ after setup.
  \item But if $n_{\rm iter}$ grows with $N/\sigma^2$, the FFT speedup is not enough.
\end{itemize}

\begin{block}{Main numerical-analysis question}
Can we reduce the iteration count without destroying the fast Fourier structure?
\end{block}
\end{frame}

% ---------------------------------------------------------
\section{EigenPro: spectral deflation preconditioning}

\begin{frame}{EigenPro intuition}
Kernel matrices often have a few very large eigenvalues:
\[
\lambda_1 \ge \lambda_2 \ge \cdots \ge \lambda_N.
\]

These top modes force small stable step sizes and slow iterative convergence.

\begin{block}{EigenPro idea}
Estimate the leading $q$ eigenvectors and construct a preconditioner that flattens the top of the spectrum.
\end{block}

\vspace{0.4em}
For a Richardson-style iteration,
\[
\alpha_{t+1}=\alpha_t-\eta (I-Q)(K\alpha_t-y),
\]
where $Q$ is a rank-$q$ spectral correction from top eigenpairs.
\end{frame}

% ---------------------------------------------------------
\begin{frame}{Deflation preconditioner: the clean algebra}
Let
\[
A = U\Theta U^*, \quad \theta_1\ge \theta_2\ge\cdots\ge\theta_n>0.
\]
Let $U_q=[u_1,\dots,u_q]$ and choose $\tau=\theta_{q+1}$. Define
\[
P_q = I-U_q\operatorname{diag}\!\left(1-\frac{\tau}{\theta_i}\right)_{i=1}^q U_q^*.
\]

Then, for exact eigenvectors,
\[
P_q A u_i =
\begin{cases}
\tau u_i, & i\le q,\\
\theta_i u_i, & i>q.
\end{cases}
\]

\begin{block}{Effect}
The top $q$ outlier eigenvalues are flattened to $\theta_{q+1}$, so
\[
\kappa(P_qA)\approx \frac{\theta_{q+1}}{\theta_{\min}(A)}
\quad \text{instead of}\quad
\frac{\theta_1}{\theta_{\min}(A)}.
\]
\end{block}
\end{frame}

% ---------------------------------------------------------
\begin{frame}{EigenPro gives the preconditioning template}
\begin{columns}[T,totalwidth=\textwidth]
\begin{column}{0.48\textwidth}
\textbf{What EigenPro does well}
\begin{itemize}
  \item Reduces effective condition number.
  \item Uses only a low-rank spectral correction.
  \item Turns large top eigenvalues into a manageable plateau.
\end{itemize}
\end{column}
\begin{column}{0.48\textwidth}
\textbf{What remains challenging}
\begin{itemize}
  \item Computing top eigenpairs can be expensive.
  \item Applying kernel matrices can still be costly.
  \item Classical EigenPro is in data/function space.
\end{itemize}
\end{column}
\end{columns}

\vspace{0.8em}
\begin{alertblock}{Thesis transition}
If EFGP gives a fast frequency-space linear system, we can transfer this spectral flattening idea to that smaller structured system.
\end{alertblock}
\end{frame}

% ---------------------------------------------------------
\section{Translation-invariant kernels and EFGP}

\begin{frame}{Translation-invariant kernels}
A kernel is translation invariant if
\[
k(x,x')=\psi(x-x').
\]

\begin{itemize}
  \item Common in spatial statistics, time series, and low-dimensional regression.
  \item Encodes a stationary prior: correlations depend on displacement, not absolute location.
  \item Examples: squared exponential, Mat\'ern, rational quadratic.
\end{itemize}

\begin{block}{Why prefer them here?}
Translation invariance exposes Fourier structure. Fourier structure exposes fast FFT/NUFFT algorithms.
\end{block}
\end{frame}

% ---------------------------------------------------------
\begin{frame}{Bochner's theorem: kernels become spectra}
For continuous positive definite translation-invariant kernels,
\[
k(t)=\int_{\R^d} e^{2\pi i \omega^\top t}\, d\mu(\omega).
\]
If a spectral density exists,
\[
k(t)=\int_{\R^d} \widehat{k}(\omega)e^{2\pi i \omega^\top t}\,d\omega,
\quad \widehat{k}(\omega)\ge 0.
\]

\begin{block}{Practical meaning}
The kernel is a weighted superposition of Fourier modes.
Smooth kernels have rapidly decaying spectra; rough kernels have heavier spectral tails.
\end{block}
\end{frame}

% ---------------------------------------------------------
\begin{frame}{EFGP: equispaced Fourier approximation}
Choose an equispaced frequency grid
\[
J_m=\{-m,\ldots,m\}^d,\quad M=(2m+1)^d.
\]
Approximate the inverse Fourier integral by a quadrature rule:
\[
k(x-x')\approx \tilde k(x-x')=
\sum_{j\in J_m} h^d\widehat{k}(hj)e^{2\pi i h j^\top(x-x')}.
\]
This is a feature factorization:
\[
\tilde k(x,x')=\sum_{j\in J_m}\phi_j(x)\overline{\phi_j(x')},
\quad
\phi_j(x)=\sqrt{h^d\widehat{k}(hj)}e^{2\pi i h j^\top x}.
\]

\begin{block}{Result}
The dense kernel matrix is approximated as $K\approx\Phi\Phi^*$ with $\Phi\in\C^{N\times M}$.
\end{block}
\end{frame}

% ---------------------------------------------------------
\begin{frame}{Weight-space system}
Instead of solving the $N\times N$ function-space system,
\[
(K+\sig I)\alpha=y,
\]
EFGP solves the $M\times M$ weight-space system
\[
(\Phi^*\Phi+\sig I)\beta = \Phi^*y.
\]
Then predictions are
\[
\hat f(x)=\sum_{j\in J_m}\beta_j\phi_j(x).
\]

\begin{block}{Key advantage}
After Fourier discretization, prediction and linear algebra depend on $M$, not directly on $N$.
\end{block}
\end{frame}

% ---------------------------------------------------------
\begin{frame}{Where FFT and NUFFT enter}
Write
\[
\Phi = F D,
\quad
F_{n,j}=e^{2\pi i h j^\top x_n},
\quad
D_{j,j}=\sqrt{h^d\widehat{k}(hj)}.
\]
Then
\[
\Phi^*\Phi=D(F^*F)D.
\]
Because frequencies are equispaced,
\[
(F^*F)_{j,\ell}=\sum_{n=1}^N e^{2\pi i h(j-\ell)^\top x_n},
\]
which depends only on $j-\ell$.

\begin{block}{Consequences}
\begin{itemize}
  \item Build Toeplitz coefficients by type-1 NUFFT: $O(N+M\log M)$.
  \item Apply $\Phi^*\Phi$ by padded FFT convolution: $O(M\log M)$ per iteration.
\end{itemize}
\end{block}
\end{frame}

% ---------------------------------------------------------
\begin{frame}{EFGP complexity summary}
\begin{center}
\begin{tabular}{lll}
\toprule
Step & Operation & Cost \\
\midrule
Precompute Toeplitz vector & NUFFT over data & $O(N+M\log M)$ \\
Build right-hand side $\Phi^*y$ & NUFFT over data & $O(N+M\log M)$ \\
One CG matvec & FFT convolution & $O(M\log M)$ \\
Predict at $q$ points & type-2 NUFFT & $O(q+M\log M)$ \\
\bottomrule
\end{tabular}
\end{center}

\vspace{0.8em}
\begin{alertblock}{Remaining issue}
The weight-space system can still be poorly conditioned, with iteration counts growing when $N/\sigma^2$ is large.
\end{alertblock}
\end{frame}

% ---------------------------------------------------------
\begin{frame}{Why preconditioning EFGP is natural}
EFGP changes the bottleneck:
\[
\text{from dense data-space kernel algebra}
\quad\longrightarrow\quad
\text{structured frequency-space linear algebra}.
\]

\begin{block}{Opportunity}
Use fast EFGP matvecs to estimate the leading spectrum of
\[
A_{\rm WS}=\Phi^*\Phi+\sig I,
\]
and apply a low-rank deflation preconditioner inside CG.
\end{block}

\begin{alertblock}{Goal}
Keep the data pass $O(N+M\log M)$, but reduce the number of FFT-CG iterations.
\end{alertblock}
\end{frame}

% ---------------------------------------------------------
\section{Proposed method}

\begin{frame}{Core idea}
\begin{block}{Frequency-space EigenPro}
Apply EigenPro-style spectral deflation to the EFGP weight-space matrix
\[
A = \Phi^*\Phi+\sig I\in\C^{M\times M}.
\]
\end{block}

\begin{enumerate}
  \item Use NUFFT/FFT to implement fast $z\mapsto Az$.
  \item Estimate top $q+1$ eigenpairs of $A$ by Lanczos or randomized subspace iteration.
  \item Build a low-rank spectral preconditioner $P_q$.
  \item Run preconditioned CG on $A\beta=b$, $b=\Phi^*y$.
\end{enumerate}

\begin{block}{Expected benefit}
Large spectral outliers are removed before CG, so residual convergence is controlled by the remaining spectrum.
\end{block}
\end{frame}

% ---------------------------------------------------------
\begin{frame}{Algorithm: preconditioned EFGP}
\small
\begin{enumerate}
  \item \textbf{Choose Fourier grid:} select $h,m$ from kernel family, tolerance, and dimension; set $M=(2m+1)^d$.
  \item \textbf{NUFFT setup:} compute Toeplitz coefficients for $F^*F$ and RHS $b=\Phi^*y$.
  \item \textbf{Fast operator:} define $A(z)=D(F^*F)(Dz)+\sig z$ using padded FFT convolution.
  \item \textbf{Spectral setup:} estimate top eigenpairs
  \[
  A u_i \approx \theta_i u_i,
  \quad i=1,\dots,q+1.
  \]
  \item \textbf{Preconditioner:}
  \[
  P_q z=z-U_q\operatorname{diag}\!\left(1-\frac{\theta_{q+1}}{\theta_i}\right)U_q^*z.
  \]
  \item \textbf{Solve and predict:} run PCG for $A\beta=b$; evaluate $\hat f(x)$ by type-2 NUFFT.
\end{enumerate}
\end{frame}

% ---------------------------------------------------------
\begin{frame}{Complexity of the proposed solver}
Let $T_A=O(M\log M)$ be the EFGP matvec cost.

\begin{center}
\begin{tabular}{lll}
\toprule
Part & Cost & Comment \\
\midrule
EFGP setup & $O(N+M\log M)$ & one pass through data \\
Eigen setup & $O(rT_A)$ & $r$ Lanczos/randomized matvecs \\
Preconditioner apply & $O(qM)$ & low-rank projection \\
PCG solve & $n_{\rm pcg}(q)(T_A+O(qM))$ & fewer iterations if spectrum is deflated \\
Prediction & $O(q_{\rm test}+M\log M)$ & type-2 NUFFT \\
\bottomrule
\end{tabular}
\end{center}

\begin{block}{Optimization problem}
Choose $q$ to minimize
\[
T(q)=T_{\rm setup}(q)+n_{\rm pcg}(q)\bigl(T_A+O(qM)\bigr).
\]
\end{block}
\end{frame}

% ---------------------------------------------------------
\begin{frame}{Why the preconditioner should help}
If eigenpairs are accurate and $\theta_1\ge\cdots\ge\theta_M$,
\[
\kappa(A)=\frac{\theta_1}{\theta_M},
\quad
\kappa(P_qA)\approx\frac{\theta_{q+1}}{\theta_M}.
\]

CG iteration count changes from roughly
\[
O\!\left(\sqrt{\frac{\theta_1}{\theta_M}}\log\frac{1}{\varepsilon}\right)
\quad\text{to}\quad
O\!\left(\sqrt{\frac{\theta_{q+1}}{\theta_M}}\log\frac{1}{\varepsilon}\right).
\]

\begin{block}{The best case}
A small number of large spectral outliers dominate the condition number; removing them gives a large iteration reduction.
\end{block}
\end{frame}

% ---------------------------------------------------------
\begin{frame}{Top-$q$ selection: the central trade-off}
\begin{columns}[T,totalwidth=\textwidth]
\begin{column}{0.48\textwidth}
\textbf{Small $q$}
\begin{itemize}
  \item Cheap eigen setup.
  \item Cheap preconditioner application.
  \item But large spectral modes remain.
  \item More CG iterations.
\end{itemize}
\end{column}
\begin{column}{0.48\textwidth}
\textbf{Large $q$}
\begin{itemize}
  \item Better spectral flattening.
  \item Fewer CG iterations.
  \item More eigenpair estimation cost.
  \item Eigenvectors/eigenvalues may be harder to estimate accurately.
\end{itemize}
\end{column}
\end{columns}

\vspace{0.8em}
\begin{alertblock}{Practical rule}
Pick $q$ near the spectral elbow: increase $q$ only while the saved CG time is larger than the spectral setup and projection cost.
\end{alertblock}
\end{frame}

% ---------------------------------------------------------
\section{Experiments}

\begin{frame}{Experimental questions}
\begin{enumerate}
  \item \textbf{Conditioning:} how does $\kappa(A)$ scale with $N$, $\sigma^2$, kernel lengthscale, and Fourier grid size?
  \item \textbf{Iteration reduction:} how much does deflation reduce CG/PCG iterations?
  \item \textbf{Runtime:} when does the eigen setup pay for itself?
  \item \textbf{Top-$q$ balance:} where is the best $q$ for total wall-clock time?
  \item \textbf{Accuracy:} does preconditioning preserve posterior mean / KRR prediction accuracy?
\end{enumerate}
\end{frame}

% ---------------------------------------------------------
\begin{frame}{Recommended experiment matrix}
\small
\begin{center}
\begin{tabular}{>{\raggedright}p{0.22\textwidth}p{0.68\textwidth}}
\toprule
Factor & Values to test \\
\midrule
Kernel & Squared exponential; Mat\'ern $\nu\in\{1/2,3/2,5/2\}$ \\
Dimension & $d=1,2,3$ where EFGP is most suitable \\
Data size & $N=10^5,10^6,10^7$ if available \\
Noise / regularization & several $\sigma^2$ or $\lambda$ values; especially small-noise regimes \\
Preconditioner size & $q\in\{0,8,16,32,64,128,256\}$ \\
Metrics & iterations, setup time, solve time, total time, residual, RMSE/EEPM \\
\bottomrule
\end{tabular}
\end{center}

\begin{block}{Baselines}
Unpreconditioned EFGP-CG, diagonal/Jacobi preconditioning, and the proposed frequency-space deflation.
\end{block}
\end{frame}

% ---------------------------------------------------------
\begin{frame}{Results slide template: runtime and iterations}
\small
Replace the bracketed entries with your measured numbers.
\begin{center}
\begin{tabular}{rrrrrr}
\toprule
$q$ & eig. setup & PCG iters & solve time & total time & rel. residual \\
\midrule
0   & --      & [ ] & [ ] & [ ] & [ ] \\
8   & [ ]     & [ ] & [ ] & [ ] & [ ] \\
16  & [ ]     & [ ] & [ ] & [ ] & [ ] \\
32  & [ ]     & [ ] & [ ] & [ ] & [ ] \\
64  & [ ]     & [ ] & [ ] & [ ] & [ ] \\
128 & [ ]     & [ ] & [ ] & [ ] & [ ] \\
\bottomrule
\end{tabular}
\end{center}

\begin{block}{Message to emphasize}
The optimal $q$ is not necessarily the largest $q$: total time balances spectral setup, preconditioner application, and iteration savings.
\end{block}
\end{frame}

% ---------------------------------------------------------
\begin{frame}{Results slide template: accuracy}
\small
Preconditioning should change convergence speed, not the target solution.
\begin{center}
\begin{tabular}{rrrrr}
\toprule
Method & $q$ & RMSE & EEPM / rel. error & final residual \\
\midrule
EFGP-CG & 0 & [ ] & [ ] & [ ] \\
Jacobi-PCG & -- & [ ] & [ ] & [ ] \\
Deflated PCG & [best] & [ ] & [ ] & [ ] \\
\bottomrule
\end{tabular}
\end{center}

\begin{alertblock}{Expected conclusion}
At matched residual tolerance, all solvers should produce nearly identical predictions; the gain should appear in iteration count and wall-clock time.
\end{alertblock}
\end{frame}

% ---------------------------------------------------------
\begin{frame}{Ablation: when does deflation help most?}
\begin{itemize}
  \item \textbf{Low noise / weak regularization}: smallest eigenvalue is close to $\sigma^2$ or $\lambda$, so conditioning worsens.
  \item \textbf{Large $N$}: largest eigenvalues grow with data scale.
  \item \textbf{Smooth kernels}: fewer Fourier modes may be needed, making eigen setup cheaper.
  \item \textbf{Spectral outliers}: clear eigengap makes top-$q$ deflation effective and stable.
\end{itemize}

\begin{block}{Failure mode}
If the spectrum decays slowly with no elbow, large $q$ may be required, and preconditioning may not amortize its setup cost.
\end{block}
\end{frame}

% ---------------------------------------------------------
\section{Summary}

\begin{frame}{Contributions to state clearly}
\begin{enumerate}
  \item Formulate EFGP's weight-space solve as a natural target for spectral deflation.
  \item Design a frequency-space preconditioner that preserves FFT/NUFFT structure.
  \item Analyze the expected condition-number improvement from top-$q$ flattening.
  \item Benchmark the top-$q$ trade-off: setup cost vs. iteration reduction.
  \item Identify regimes where preconditioning is most valuable: large $N/\sigma^2$, low dimension, translation-invariant kernels.
\end{enumerate}
\end{frame}

% ---------------------------------------------------------
\begin{frame}{Limitations and future work}
\begin{itemize}
  \item \textbf{Dimensionality}: EFGP uses $M=(2m+1)^d$ modes, so high dimension remains difficult.
  \item \textbf{Eigenpair quality}: approximate top eigenvectors must be accurate enough to keep the preconditioner stable.
  \item \textbf{Choosing $q$}: current rule is empirical; a more automatic spectral-elbow criterion is desirable.
  \item \textbf{Theory}: need sharper bounds relating the deflated weight-space spectrum to data distribution and kernel parameters.
  \item \textbf{Extensions}: posterior variance, hyperparameter learning, GPU implementation, and randomized eigensolvers.
\end{itemize}
\end{frame}

% ---------------------------------------------------------
\begin{frame}{Takeaway}
\begin{center}
\Large
KRR is not obsolete; its bottleneck is linear algebra.
\end{center}

\vspace{1em}
\begin{block}{Main thesis insight}
EFGP makes low-dimensional translation-invariant KRR fast per iteration; EigenPro-style spectral deflation can make it fast in iteration count.
\end{block}

\vspace{0.5em}
\begin{center}
\Large
\textbf{Fast Fourier structure + spectral preconditioning}
\end{center}
\end{frame}

% ---------------------------------------------------------
\begin{frame}[allowframebreaks]{References}
\small
\begin{thebibliography}{9}
\bibitem{greengard2023}
P. Greengard, M. Rachh, and A. Barnett.
\newblock Equispaced Fourier representations for efficient Gaussian process regression from a billion data points.
\newblock 2023.

\bibitem{barnett2023}
A. Barnett, P. Greengard, and M. Rachh.
\newblock Uniform approximation of common Gaussian process kernels using equispaced Fourier grids.
\newblock 2023.

\bibitem{eigenpro3}
A. Abedsoltan, M. Belkin, and P. Pandit.
\newblock Toward Large Kernel Models.
\newblock ICML 2023.

\bibitem{doumeche2025}
N. Doum\`eche, F. Bach, G. Biau, and C. Boyer.
\newblock Fast Kernel Methods: Sobolev, Physics-Informed, and Additive Models.
\newblock 2025.

\bibitem{bach2024}
F. Bach.
\newblock Learning Theory from First Principles.
\newblock MIT Press, 2024.
\end{thebibliography}
\end{frame}

% ---------------------------------------------------------
\appendix
\section{Backup slides}

\begin{frame}{Backup: function space vs. weight space equivalence}
If $K\approx \Phi\Phi^*$, then the function-space approximate solve is
\[
(\Phi\Phi^*+\sig I)\tilde\alpha=y.
\]
The weight-space solve is
\[
(\Phi^*\Phi+\sig I)\beta=\Phi^*y.
\]
They are linked by
\[
\beta=\Phi^*\tilde\alpha,
\qquad
\tilde\alpha=\frac{y-\Phi\beta}{\sig}.
\]
Therefore the posterior mean can be evaluated as
\[
\mu(x)=\sum_j \beta_j\phi_j(x).
\]
\end{frame}

\begin{frame}{Backup: kernel approximation error in EFGP}
The Fourier grid approximation error can be decomposed into
\[
\tilde{k}(x)-k(x)
= -\sum_{n\in\mathbb{Z}^d,n\ne 0} k\!\left(x+\frac{n}{h}\right)
+ h^d\sum_{j\in\mathbb{Z}^d,j\notin J_m}\widehat{k}(hj)e^{2\pi i h j^\top x}.
\]

\begin{itemize}
  \item First term: aliasing error from grid spacing $h$.
  \item Second term: truncation error from finite frequency box $J_m$.
\end{itemize}
\end{frame}

\begin{frame}{Backup: preconditioner stability checklist}
\begin{itemize}
  \item Verify eigenvalues are positive and ordered.
  \item Clip factors $1-\tau/\theta_i$ to avoid negative preconditioner weights under numerical noise.
  \item Check $\norm{U_q^*U_q-I}$ for orthogonality.
  \item Monitor PCG residual in the original system, not only the preconditioned residual.
  \item Compare final predictions against unpreconditioned CG at matched tolerance.
\end{itemize}
\end{frame}

\end{document}
